\documentclass[10pt,a4paper]{amsart}
\usepackage[margin=19mm]{geometry}
\usepackage{amsmath,amssymb,mathtools}
\usepackage[scr=rsfs,cal=euler]{mathalfa}
\usepackage{xfrac}
\usepackage[unicode,hidelinks,bookmarks=false]{hyperref}
\newcommand{\ie}{i.e.\ }
\newcommand{\cf}{cf.\ }
\newcommand{\resp}{respectively\ }
\newcommand{\Subsection}{\subsection}
\providecommand{\nobreakdash}{\nobreak\hbox{-}}
\setlength{\emergencystretch}{2em}
\multlinegap0pt
\usepackage{amssymb}
\usepackage[shortlabels]{enumitem}
\usepackage{mathtools}
\usepackage{bbm}
\newtheorem{theorem}{Theorem}
\newtheorem{lemma}{Lemma}
\newtheorem{construction}{Construction}
\title{Ekedahl-Oort Types and Newton Polygons of Abelian Covers of $\mathbf{P}^1$ Branched at Three Points}
\author{Darren Schmidt}
\date{Source: arXiv:2602.07693v1 (CC BY 4.0)}
\begin{document}
\maketitle

\noindent\emph{Source and licence.} Ekedahl-Oort Types and Newton Polygons of Abelian Covers of $\mathbf{P}^1$ Branched at Three Points. Version arXiv:2602.07693v1 (CC BY 4.0), distributed by arXiv under the Creative Commons Attribution 4.0 International licence, http://creativecommons.org/licenses/by/4.0/, as recorded in the arXiv licence metadata for this exact version and shown on the abstract page https://arxiv.org/abs/2602.07693v1. Full licence notice: \texttt{licenses/arxiv-2602.07693-CC-BY-4.0.txt}.\par\smallskip

\noindent\textit{Editorial scope.} Counted unit: the article's theorem slopes (source0.tex lines 109--112). The article's complete original proof of it is delivered with it (source0.tex lines 541--615, one \texttt{proof} environment of 75 lines, which the article places away from the statement). No mathematics is written, repaired or re-derived: the statement and the proof are the article's own lines, in the article's own notation, and no retained line is substituted or re-flowed.  Retained uncounted as the source-local context the proof needs: lines 355--358; lines 405--408; lines 421--427; lines 492--522.  Nothing was omitted from any retained span, so the package carries no omission record.  No retained line contains a citation, so the wrapper carries no bibliography and no bibliography entry.  No retained line contains a figure environment, an image-inclusion command or drawing code, so no image is delivered and the package declares no asset dependency.  Editorial formatting only, in addition: an AMS \texttt{amsart} preamble that restates the article's own declarations and macros used by the retained lines, namely the environments \texttt{theorem}, \texttt{lemma}, \texttt{construction} on separate counters, together with the standard packages those lines need.  The retained environments print in this document as Theorem 1, Lemma 1, Construction 1 in order of appearance, whereas the article prints the same items with the numbering of its own full document.  The complete native source of the article as distributed in the arXiv e-print for this exact version is bundled unchanged as \texttt{sources/194199/source0.tex}.
\begin{lemma}
\label{lemma}
    Let $\ell > 3$ be a prime. If $n$ and $r$ are positive integers such that $n \geq r$, then \begin{align}\left (\frac{-1}{\ell} \right ) q\left(\frac{\ell r}{4n}, \frac{4\ell r}{4n} \right) = \sum_{i=1}^3 q\left (\frac{\ell(in-r)}{4n}, \frac{i\ell n}{4n} \right) \label{lemma-align}. \end{align}
\end{lemma}

\begin{construction}
    \label{construction}
    Let $\ell > 3$ be a prime and $n \geq 1$ such that gcd$(n,\ell) = 1$. Pick $0 \leq k \leq n-1$ such that $k \equiv -3\ell^{-1}$ modulo $n$.  Consider the cover $\phi: X \to \mathbf{P}^1$ branched at three points with ramification type $[n\ell,n\ell,\ell,1]$ and inertia type $(1,(n-k)\ell-4,k\ell+3)$. The equation for $X$ is $y^{n\ell} = x(1-x)^{(n-k)\ell-4}$.
\end{construction}

\begin{lemma}
\label{sigtype}
    Let $\phi$ be the cover in Construction \ref{construction}, and let $f$ be the signature type of $\phi$. If $j \in \{1,\ldots,\ell n-1\}$ and $j \not \equiv 0$ modulo $\ell$, pick $i$ such that $0 \leq i \leq 6n-1$ and $j \in I(i)$. Then,
    \begin{align*}
        f(j) = \left \lceil \frac{jk+\left \lfloor \frac{i+2}{2}\right \rfloor }{n}\right \rceil - \left \lfloor \frac{jk + \left \lfloor \frac{i+1}{2}\right \rfloor + \left \lfloor \frac{i}{6} \right \rfloor }{n} \right \rfloor
    \end{align*}
\end{lemma}

\begin{lemma}
    \label{intervals}
    Let $\phi$ be the cover in Construction \ref{sigtype}, and let $f$ be the signature type of $\phi$. Let $j \in \{1,2,\ldots,\ell n-1\}$ such that $j \equiv a\ell$ modulo $n$ for some $0 \leq a \leq n-1$ and $j \not \equiv 0$ modulo $\ell$. We know that $f(j) = 1$ if and only if
    \begin{itemize}

   
    \item $0 \leq 3a \leq n-1$ and
    
    \[
        %\begin{split}
        j \in \left(0, \frac{3a\ell}{4}\right) \cup \left(a\ell, \frac{\ell(n+3a)}{4} \right) \cup \left(\frac{\ell(n+3a)}{3}, \frac{\ell(2n+3a)}{4} \right) \cup \left(\frac{\ell(2n+3a)}{3}, \frac{\ell(3n+3a)}{4} \right),
        %\end{split}
    \]
    
    \item or $n \leq 3a \leq 2n-1$ and

    \[
        %\begin{split}
        j \in \left(0, \frac{\ell(3a-n)}{4}\right) \cup \left(\frac{\ell(3a-n)}{3}, \frac{3\ell a}{4} \right) \cup  \left(a\ell, \frac{\ell(n+3a)}{4} \right) \cup  \left(\frac{\ell(n+3a)}{3}, \frac{\ell(2n+3a)}{4} \right),
       % \end{split}
    \]
    
    \item or $2n \leq 3a \leq 3n-1$ and
   
    \[
       % \begin{split}
       j \in \left(0, \frac{\ell(3a-2n)}{4}\right) \cup \left(\frac{\ell(3a-2n)}{3}, \frac{\ell(3a-n)}{4} \right) \cup  \left(\frac{\ell(3a-n)}{3}, \frac{3\ell a}{4} \right) \cup  \left(a\ell, \frac{\ell(n+3a)}{4} \right).
        %\end{split}
    \]
     \end{itemize}
\end{lemma}

\begin{theorem}
\label{slopes}
Let $\ell > 3$ be a prime and $n \geq 1$ be coprime to $\ell$. Define $g := (\ell-1)/2$ and suppose $\alpha$ is the number of quadratic residues modulo $\ell$ contained in the intervals $(0,\ell/4)$, $(\ell/3,\ell/2)$, and $(2\ell/3, 3\ell/4)$. If $p$ is a prime with order $g$ in $\left (\mathbf{Z}/\ell\mathbf{Z} \right)^\times$ such that $gcd(p,n\ell) = 1$, and the order of $p$ in $\left (\mathbf{Z}/n\mathbf{Z} \right)^\times$ is coprime to $g$, then there exists a cyclic cover $\phi: X \to \mathbf{P}^1$ branched at three points defined over $\overline{\mathbf{F}_p}$ such that $X$ has genus $ng$ and the only slopes of the the Newton polygon of $\textrm{Jac}(X)$ are $\alpha/g$ and $(g-\alpha)/g$.
\end{theorem}

\begin{proof}[Proof of Theorem \ref{slopes}]
Let $f$ be the signature type of the cover in Construction \ref{construction}. Let $j \in \{1,2,\ldots,\ell n-1\}$ such that $j \equiv a\ell$ modulo $n$ for some $0 \leq a \leq n-1$ and $j \not \equiv 0$ modulo $\ell$.

First, assume $0 \leq 3a \leq n-1$, then Lemma \ref{intervals} says that $f(j) =1$ if and only if 

\begin{align}
0 < j& < \frac{3a\ell}{4} \label{first_i} \\
a\ell < j& < \frac{\ell(n+3a)}{4} \label{second_i} \\
 \frac{(n+3a)\ell}{3}  < j& < \frac{(2n+3a)\ell}{4} \label{third_i} \\
 \frac{(2n+3a)\ell}{3}  <  j& <  \frac{(3n+3a)\ell}{4} \label{fourth_i}
\end{align}

We wish to compute the number of quadratic residues and nonresidues inside of 
\begin{align} \left(0, \frac{3a\ell}{4}\right) \cup \left(a\ell, \frac{\ell(n+3a)}{4} \right) \cup \left(\frac{\ell(n+3a)}{3}, \frac{\ell(2n+3a)}{4} \right) \cup \left(\frac{\ell(2n+3a)}{3}, \frac{\ell(3n+3a)}{4} \right)\label{first-a-interval} \end{align}

and show that it has the same number of quadratic residues and nonresidues as 
\begin{align}\left(0, \frac{\ell}{4} \right) \cup \left(\frac{\ell}{3}, \frac{\ell}{2} \right) \cup \left(\frac{2\ell}{3}, \frac{3\ell}{4} \right). \label{needed_interval} \end{align}

We note that intervals $(\ref{first-a-interval})$ and $(\ref{needed_interval})$ contain $g$ integers, which is the sum of the number of non-zero quadratic residues and nonresidues in these intervals. The quadratic excess of any of these intervals is the number of quadratic residues minus the number of quadratic nonresidues in the interval. Therefore, if the quadratic excess of these intervals is equal, then they have the same number of quadratic residues and the same number of quadratic nonresidues. The number of non-zero quadratic residues in this interval gives one slope of the Newton polygon, and the number of quadratic nonresidues gives the other slope.

As $j \equiv a\ell$ modulo $n$, we can add $(n-1)a\ell$ to interval (\ref{first_i}) to get elements congruent to $0$ modulo $n$. As we are adding a multiple of $\ell$, this doesn't change the number of quadratic residues or nonresidues in the interval. For intervals (\ref{second_i}), (\ref{third_i}), and (\ref{fourth_i}), subtract $a\ell$. Then, divide each inequality by $n$. Finally, we shift the resulting interval from interval 
(\ref{first_i}) by subtracting $(a-1)\ell$.  We now need to compute

\begin{align}
    q\left(\frac{3\ell(n-a)}{3n}, \frac{\ell(4n-a)}{4n}\right) + q\left(0, \frac{\ell(n-a)}{4n} \right) + q\left(\frac{n\ell}{3n}, \frac{\ell(2n-a)}{4n} \right) + q\left(\frac{2n\ell}{3n}, \frac{\ell(3n-a)}{4n} \right). \label{first-qe}
\end{align}

We see that

\begin{align}
    q\left(\frac{3\ell(n-a)}{3n}, \frac{\ell(4n-a)}{4n}\right) = \left (\frac{-1}{\ell} \right) q \left(\frac{a\ell}{4n}, \frac{3a\ell}{3n} \right) = \left (\frac{-1}{\ell} \right) q \left(\frac{a\ell}{4n}, \frac{4a\ell}{4n} \right) \label{interval-manip}
\end{align}

Finally, by substituting equation (\ref{interval-manip}) into Lemma \ref{lemma} with $r = a$, we get that equation (\ref{first-qe}) equals

\begin{align*}
    q\left(0, \frac{\ell}{4} \right) + q\left(\frac{\ell}{3}, \frac{\ell}{2} \right) + q\left(\frac{2\ell}{3}, \frac{3\ell}{4} \right).
\end{align*}

This concludes the case when $0 \leq 3a \leq n-1$.

Assume $n \leq 3a \leq 2n-1$, then Lemma \ref{intervals} states that if $j \equiv a\ell$ modulo $n$ and $j \not \equiv 0 $ modulo $\ell$, then $f(j) = 1$ if and only if $j$ lies in one of the following intervals:

\begin{align}
    0 < j& < \frac{\ell(3a-n)}{4} \label{2a-1}\\
    \frac{\ell(3a-n)}{3} < j& < \frac{3\ell a}{4} \label{2a-2}\\
    a\ell < j& < \frac{\ell(n+3a)}{4} \label{2a-3}\\
    \frac{\ell(n+3a)}{3} < j& < \frac{\ell(2n+3a)}{4} \label{2a-4}
\end{align}

Note that intervals (\ref{2a-3}) and (\ref{2a-4}) are the same as intervals (\ref{second_i}) and (\ref{third_i}), respectively, so we only need to prove that the quadratic excess of the intervals (\ref{2a-1}) and (\ref{2a-2}) equals the quadratic excess of intervals (\ref{first_i}) and (\ref{fourth_i}). We add $a\ell(n-1)$ to both (\ref{2a-1}) and (\ref{2a-2}), divide by $n$, and then subtract $\ell(a-1)$. We now count the quadratic excess

\begin{align*}
   q&\left(\frac{3\ell(n-a)}{3n}, \frac{\ell(3n-a)}{4n}\right) + q\left(\frac{2n\ell}{3n}, \frac{\ell(4n-a)}{4n}\right)  \\ =  q&\left(\frac{3\ell(n-a)}{3n}, \frac{\ell(4n-a)}{4n} \right) - q\left(\frac{\ell(3n-a)}{4n} \frac{\ell(4n-a)}{4n} \right) \\ + \ q&\left(\frac{2n\ell}{3n}, \frac{\ell(3n-a)}{4n} \right) + q\left(\frac{\ell(3n-a)}{4n},\frac{\ell(4n-a)}{4n} \right)
\end{align*}

Therefore, we get the same quadratic excess as in equation (\ref{first-qe}).

For the final case, assume that $2n \leq 3a \leq 3n-1$. Lemma \ref{intervals} state that when $j \equiv a\ell$ modulo $n$, $j \not \equiv 0$ modulo $\ell$, and $f(j) = 1$:

\begin{align}
    0 < j& < \frac{\ell(3a-2n)}{4} \label{3a-1} \\
    \frac{\ell(3a-2n)}{3} < j& < \frac{\ell(3a-n)}{4} \label{3a-2} \\
    \frac{\ell(3a-n)}{3} < j & < \frac{3\ell a}{4} \label{3a-3} \\
    a\ell < j& < \frac{\ell(n+3a)}{4} \label{3a-4}
\end{align}

As intervals (\ref{3a-3}) and (\ref{3a-4}) match intervals (\ref{2a-2}) and (\ref{2a-3}), respectively, we only have to focus on intervals (\ref{3a-1}) and (\ref{3a-2}). We get the quadratic excess

\begin{align*}
    q\left(\frac{3\ell(n-a)}{3n}, \frac{\ell(2n-a)}{4n} \right) + q\left(\frac{n\ell}{3n}, \frac{\ell(3n-a)}{4n} \right)
\end{align*}

and this equals the quadratic excess coming from intervals (\ref{2a-1}) and (\ref{2a-4}).
\end{proof}

\end{document}
